% =========================================================
% frontmatter/portfolio_context.tex
% =========================================================

\chapter*{Portfolio Context and Evidence}
\thispagestyle{fancy}
\addcontentsline{toc}{chapter}{Portfolio Context and Evidence}

\noindent
\textbf{Project Context.}
This work originated from a graduate-level antenna engineering project
completed during Summer 2026. The present document is an independently
prepared public portfolio edition intended to demonstrate analytical
antenna design, full-wave electromagnetic simulation, array analysis,
beam steering, engineering interpretation, and technical documentation.

\vspace{0.5em}

\noindent
\textbf{Public Portfolio Edition.}
The original academic report has been reorganized and edited for
professional and educational use. Course-specific grading material,
student-identification information, institutional branding, and
administrative content have been removed. This document does not
represent an official publication or endorsement by any university,
instructor, software vendor, or other institution.

\vspace{0.5em}

\noindent
\textbf{Evidence Classification.}
The results presented in this report are based on analytical calculations
and CST Studio Suite full-wave simulations. No antenna prototype was
fabricated, and no experimental measurement campaign was conducted.
Accordingly, the reported resonant frequencies, scattering parameters,
radiation characteristics, efficiencies, directivities, and
beam-steering results are classified as analytical or simulated evidence,
not measured performance.

\vspace{0.5em}

\noindent
\textbf{Simulation Scope.}
The baseline models were developed using CST Studio Suite Learning
Edition. Because of the available mesh-cell limitation, PEC
metallization was used for the successfully tuned baseline simulations
while dielectric substrate loss was retained. A supplementary
fixed-mesh finite-conductivity copper simulation was performed for the
single patch to evaluate conductor-loss sensitivity. The numerical
limitations and their implications are discussed explicitly in the
engineering-reliability section of this report.

\vfill

\begin{center}
    \textit{Prepared as part of the Hossein Molhem Engineering Portfolio.}
\end{center}
